\chapter{Aritmètica ordinal}\label{cap:aritmetica-ordinal}

\begin{objectius}
\apartat{Prerequisits} Del capítol~\ref{cap:ordinals}: tipus ordinals, el seu ordre, successor, límit i suprem acotat. Del capítol~\ref{cap:bons-ordres}: rigidesa, segments i comparació. Del capítol~\ref{cap:recursio}: recursió transfinita amb codomini fix. Del capítol~\ref{cap:zorn}: el conjunt de Hartogs com a cota. Dels capítols~\ref{cap:valors} i~\ref{cap:parts}: sumes i productes de conjunts.
\apartat{Objectius} En acabar el capítol, el lector hauria de poder:
\begin{itemize}
  \item definir la suma i el producte ordinals sobre representants (concatenació; \(\beta\) còpies d'\(\alpha\) amb l'ordre antilexicogràfic), demostrar que són bons ordres ben definits i associatius;
  \item calcular \(1+\omega\), \(\omega+1\), \(2\cdot\omega\), \(\omega\cdot2\), i demostrar la no-commutativitat i la distributivitat només per l'esquerra;
  \item definir l'exponenciació per funcions de suport finit, deduir-ne les equacions recursives per a base no nul·la, i explicar per què la base \(0\) es tracta a part;
  \item demostrar que una immersió d'ordre entre bons ordres compara els tipus, i deduir-ne la monotonia de les tres operacions;
  \item distingir l'aritmètica ordinal (recorda l'ordre) de la cardinal (l'oblida), amb \(2^{\omega}=\omega\) com a exemple clau.
\end{itemize}
\end{objectius}

Els tipus ordinals mesuren la forma d'un bon ordre; ara els operem. La suma,
el producte i l'exponenciació ordinals es defineixen ---seguint la disciplina
del capítol anterior--- \emph{sobre representants}, amb construccions
concretes de bons ordres, i després es comprova que el resultat no depèn del
representant. Dues coses fan aquesta aritmètica diferent de la dels naturals,
i val la pena anunciar-les d'entrada. Primer, \emph{no és commutativa}: l'ordre
en què es concatenen o s'iteren els bons ordres importa, i \(1+\omega\) no és
\(\omega+1\). Segon, recorda \emph{l'ordre} i no la \emph{mida}: dos bons
ordres poden ser equipotents ---tenir la mateixa cardinalitat--- i tenir tipus
ordinals ben diferents; l'aritmètica ordinal distingeix el que la cardinal
(capítol següent) oblidarà. La suma i el producte són operacions binàries de
Nivell~I, construïdes directament sobre representants. L'exponenciació es
defineix per recursió transfinita (capítol~\ref{cap:recursio}), i la seva
clàusula límit ---un suprem--- és el primer lloc on cal vigilar l'acotació:
mentre la base i l'exponent visquin dins un \(\Hart(A)\) prou gran, tot és de
Nivell~I; una teoria completament general demanaria SR.

\section{Suma ordinal}\label{sec:suma-ordinal}

\begin{definicio}[Suma ordinal]\index{Suma ordinal}\label{def:suma-ordinal}
Siguin \(\alpha=\operatorname{tipus}(A)\) i \(\beta=\operatorname{tipus}(B)\)
tipus ordinals amb representants \((A,<_A)\) i \((B,<_B)\). La \emph{suma
ordinal} \(\alpha+\beta\) és el tipus del bon ordre sobre la suma disjunta
\(A+B\) (definició~\ref{def:suma-disjunta}) en què \emph{tots els elements
d'\(A\) precedeixen tots els de \(B\)}, i dins de cada sumand es manté l'ordre
original:
\[
x<y
\ \Longleftrightarrow\
\begin{cases}
x<_A y & \text{si } x,y\in A,\\
x<_B y & \text{si } x,y\in B,\\
\text{cert} & \text{si } x\in A,\ y\in B,\\
\text{fals} & \text{si } x\in B,\ y\in A.
\end{cases}
\]
\end{definicio}

\begin{proposicio}[La suma és un bon ordre, invariant]\label{prop:suma-bon-ordre}
L'ordre de la definició~\ref{def:suma-ordinal} és un bon ordre sobre \(A+B\),
i el seu tipus no depèn dels representants: si \(A\cong A'\) i \(B\cong B'\)
com a ordres, aleshores \(A+B\cong A'+B'\). Per tant \(\alpha+\beta\) està ben
definida.\par
\emph{Axiomes usats:} suma disjunta (capítol~\ref{cap:valors}), bon ordre
(capítol~\ref{cap:bons-ordres}), unicitat de l'isomorfisme d'ordre
(teorema~\ref{teo:unicitat-iso}); cap elecció.
\end{proposicio}

\begin{prova}
\emph{És un ordre estricte total}: la irreflexivitat i la transitivitat es
comproven per casos segons els sumands (els casos mixtos són consistents
perquè \enquote{\(A\) abans de \(B\)} és una decisió global), i la totalitat
val perquè dos elements qualssevol o són al mateix sumand ---on \(<_A\) o
\(<_B\) els compara--- o a sumands diferents ---on la clàusula mixta els
compara. \emph{És ben ordenat}: sigui \(P\hookrightarrow A+B\) una part no
inicial. Si \(P\cap A\) és no inicial, el seu mínim (pel bon ordre d'\(A\)) és
el mínim de \(P\), perquè tot element d'\(A\) precedeix tot element de \(B\).
Si \(P\cap A\) és inicial, aleshores \(P\subseteq B\) i el mínim de \(P\) és el
seu mínim pel bon ordre de \(B\). \emph{Invariància}: dos isomorfismes d'ordre
\(f\colon A\cong A'\) i \(g\colon B\cong B'\) donen \(f+g\colon A+B\cong A'+B'\),
que conserva l'ordre perquè respecta cada sumand i la relació
\enquote{\(A\) abans de \(B\)}.
\end{prova}

\begin{proposicio}[Associativitat de la suma]\label{prop:suma-associativa}
Per a tipus ordinals qualssevol, \((\alpha+\beta)+\gamma=\alpha+(\beta+\gamma)\).
\end{proposicio}

\begin{prova}
Els dos costats són el tipus del bon ordre sobre \(A+B+C\) en què tot \(A\)
precedeix tot \(B\), que precedeix tot \(C\), amb els ordres originals dins de
cada sumand. L'associativitat de la suma disjunta de conjunts
(capítol~\ref{cap:valors}) dona un isomorfisme \((A+B)+C\cong A+(B+C)\) que
conserva aquest ordre; per la unicitat de l'isomorfisme d'ordre
(teorema~\ref{teo:unicitat-iso}), els tipus coincideixen.
\end{prova}

\begin{proposicio}[La suma no és commutativa]\label{prop:suma-no-commutativa}
En general \(\alpha+\beta\neq\beta+\alpha\). Concretament,
\[
1+\omega=\omega,
\qquad\text{però}\qquad
\omega+1\neq\omega .
\]
\end{proposicio}

\begin{prova}
\emph{\(1+\omega=\omega\)}: \(1+\omega\) és el tipus d'un punt seguit d'una
còpia de \(\N\); l'aplicació que envia el punt afegit a \(0\) i el natural
\(n\) a \(n+1\) és un isomorfisme d'ordre amb \(\N\) ---simplement desplaça la
numeració---, de manera que \(1+\omega=\operatorname{tipus}(\N)=\omega\).
\emph{\(\omega+1\neq\omega\)}: \(\omega+1\) és el tipus d'una còpia de \(\N\)
seguida d'un punt \(\ast\) col·locat \emph{després} de tots els naturals.
Aquest bon ordre té un element \emph{màxim} (el punt \(\ast\)), mentre que
\(\N\) no en té cap; com que tenir màxim és invariant per isomorfisme d'ordre,
\(\omega+1\not\cong\omega\), és a dir \(\omega+1\neq\omega\). De fet
\(\omega+1\) és el successor de \(\omega\) (definició~\ref{def:zero-successor}),
estrictament més gran que \(\omega\), mentre que \(1+\omega=\omega\): la suma
per l'esquerra d'un punt \enquote{s'absorbeix} i per la dreta no.
\end{prova}

\begin{observacio}[Lectura clàssica: concatenar ordres]\label{obs:suma-classica}
La suma ordinal és la \emph{concatenació} de bons ordres: \(\alpha+\beta\) és
\enquote{primer tot \(\alpha\), després tot \(\beta\)}. La no-commutativitat
és intuïtiva des d'aquest punt de vista: afegir un element \emph{abans} d'una
successió infinita ascendent no canvia la forma (es pot reabsorbir desplaçant
els índexs), però afegir-lo \emph{després} crea un element màxim que abans no
hi era, i canvia la forma. És el primer exemple, i el més net, de la
diferència entre ordre i mida: \(1+\omega\) i \(\omega+1\) són equipotents
---tots dos numerables--- però tenen tipus ordinals diferents.
\end{observacio}

\section{Producte ordinal}\label{sec:producte-ordinal}

\begin{definicio}[Producte ordinal]\index{Producte ordinal}\label{def:producte-ordinal}
Siguin \(\alpha=\operatorname{tipus}(A)\) i \(\beta=\operatorname{tipus}(B)\).
El \emph{producte ordinal} \(\alpha\cdot\beta\) és el tipus del bon ordre
sobre el producte \(A\times B\) ordenat \emph{primer per la coordenada de
\(B\) i després per la d'\(A\)} ---l'ordre \emph{antilexicogràfic}:
\[
(a,b)<(a',b')
\quad\Longleftrightarrow\quad
b<_B b'\ \ \text{o bé}\ \ (b=b'\ \text{i}\ a<_A a').
\]
Intuïtivament, \(\alpha\cdot\beta\) són \(\beta\) còpies d'\(\alpha\)
col·locades una rere l'altra: per a cada \(b\in B\) hi ha una còpia sencera
d'\(A\), i les còpies se succeeixen segons l'ordre de \(B\). Fixem aquesta
convenció ---la coordenada de \(B\) mana--- abans de qualsevol càlcul.
\end{definicio}

\begin{proposicio}[El producte és un bon ordre, invariant i associatiu]\label{prop:producte-bon-ordre}
L'ordre antilexicogràfic de la definició~\ref{def:producte-ordinal} és un bon
ordre sobre \(A\times B\); el seu tipus no depèn dels representants; i el
producte és associatiu, \((\alpha\cdot\beta)\cdot\gamma=\alpha\cdot(\beta\cdot\gamma)\).\par
\emph{Axiomes usats:} producte (capítol~\ref{cap:parts}), bon ordre, unicitat
de l'isomorfisme d'ordre (teorema~\ref{teo:unicitat-iso}); cap elecció.
\end{proposicio}

\begin{prova}
\emph{Bon ordre}: és un ordre estricte total (l'antilexicogràfic d'ordres
totals ho és). Per al bon ordre, sigui \(P\hookrightarrow A\times B\) no
inicial; la projecció \(\pi_B(P)\hookrightarrow B\) és no inicial i té un mínim
\(b_0\); dins la fibra \(\{(a,b_0)\in P\}\), la projecció a \(A\) és no
inicial i té un mínim \(a_0\); aleshores \((a_0,b_0)\) és el mínim de \(P\),
perquè cap element amb segona coordenada \(>b_0\) no el precedeix, i dins
\(b_0\) mana \(a_0\). \emph{Invariància}: \(f\times g\) transporta l'ordre.
\emph{Associativitat}: tots dos costats són el tipus de \(A\times B\times C\)
amb l'ordre \enquote{la coordenada de \(C\) primer, després la de \(B\),
després la d'\(A\)}; l'isomorfisme canònic \((A\times B)\times C\cong A\times(B\times C)\)
el conserva, i la unicitat (teorema~\ref{teo:unicitat-iso}) tanca.
\end{prova}

\begin{proposicio}[El producte no és commutatiu; distributivitat per l'esquerra]\label{prop:producte-propietats}
En general \(\alpha\cdot\beta\neq\beta\cdot\alpha\); concretament
\[
2\cdot\omega=\omega,
\qquad\text{però}\qquad
\omega\cdot2=\omega+\omega\neq\omega .
\]
A més val la distributivitat \emph{per l'esquerra}:
\(\alpha\cdot(\beta+\gamma)=\alpha\cdot\beta+\alpha\cdot\gamma\); en canvi la
distributivitat per la dreta \emph{falla} en general.
\end{proposicio}

\begin{prova}
\emph{\(2\cdot\omega=\omega\)}: són \(\omega\) còpies del bon ordre de dos
elements, és a dir la successió \(0<1<2<3<\cdots\) agrupada de dos en dos
\((0,1),(2,3),\dots\); reindexant, és isomorfa a \(\N\), d'on
\(2\cdot\omega=\omega\). \emph{\(\omega\cdot2=\omega+\omega\)}: són \(2\)
còpies de \(\N\), l'una després de l'altra ---exactament la concatenació
\(\omega+\omega\), que té un element que no és el mínim i no té predecessor
immediat (el primer de la segona còpia); \(\N\) no en té cap de tal, perquè
tot punt no nul és un successor (lema~\ref{lema:zero-o-successor}), i la
propietat és invariant per isomorfisme d'ordre. Per tant
\(\omega+\omega\not\cong\omega\). Per
tant \(\omega\cdot2\neq2\cdot\omega\). \emph{Distributivitat esquerra}:
\(\alpha\cdot(\beta+\gamma)\) és \((\beta+\gamma)\) còpies d'\(\alpha\), és a
dir \(\beta\) còpies seguides de \(\gamma\) còpies, que és
\(\alpha\cdot\beta+\alpha\cdot\gamma\); l'isomorfisme d'ordre és el canònic
\(A\times(B+C)\cong(A\times B)+(A\times C)\), que conserva l'antilexicogràfic.
\emph{La dreta falla}: \((\ 1+1)\cdot\omega=2\cdot\omega=\omega\), però
\(1\cdot\omega+1\cdot\omega=\omega+\omega\neq\omega\).
\end{prova}

\begin{observacio}[Lectura clàssica: substituir cada punt per una còpia]\label{obs:producte-classic}
El producte \(\alpha\cdot\beta\) s'obté prenent el bon ordre \(\beta\) i
substituint \emph{cada} un dels seus punts per una còpia sencera d'\(\alpha\).
La convenció \enquote{la coordenada de \(\beta\) mana} és la que fa que això
sigui \(\beta\) còpies d'\(\alpha\) i no a l'inrevés; amb la convenció
oposada s'obtindria \(\beta\cdot\alpha\), i com que el producte no és
commutatiu, cal fixar-la. La distributivitat només val per l'esquerra
perquè substituir cada punt d'una concatenació \(\beta+\gamma\) per una còpia
d'\(\alpha\) dona la concatenació de les substitucions; en canvi
\enquote{repartir} una còpia d'\(\alpha\) sobre una suma a la dreta trenca
l'estructura.
\end{observacio}

\section{Exponenciació ordinal}\label{sec:exponenciacio-ordinal}

L'exponenciació no es defineix per una construcció directa sinó \emph{per
recursió} sobre l'exponent, amb un suprem als límits. És on la disciplina de
l'acotació esdevé imprescindible.

Per evitar una recursió sobre \enquote{tots els tipus}, que demanaria un
codomini fix i una cota per a cada suprem, donem una construcció
\emph{directa} sobre representants, com hem fet amb la suma i el producte,
i en deduïm les equacions recursives com a teorema.

\begin{definicio}[Funcions de suport finit]\index{Funcions de suport finit}\label{def:suport-finit}
Siguin \((A,<_A)\) i \((B,<_B)\) bons ordres, amb \(A\) no buit i mínim
\(0_A\). Una aplicació \(f\colon B\to A\) té \emph{suport finit} si la part
\(\operatorname{supp}(f):=\{b:f(b)\neq0_A\}\hookrightarrow B\) és finita
(capítol~\ref{cap:finits}). Les aplicacions de suport finit formen una part
\[
E(B,A)\ \hookrightarrow\ A^{B},
\]
tallada per la condició interna \enquote{la part classificada per
\(\neg\circ\chi_{0_A}\circ\operatorname{ev}\) és finita}. L'ordenem
\emph{antilexicogràficament}: \(f<g\) si \(f\neq g\) i, al \emph{màxim}
punt \(b\) on difereixen, \(f(b)<_Ag(b)\). El màxim existeix perquè els
punts on difereixen formen una part finita no buida
(\(\subseteq\operatorname{supp}f\cup\operatorname{supp}g\)) d'un ordre
total.
\end{definicio}

\begin{definicio}[Exponenciació ordinal]\index{Exponenciació ordinal}\label{def:exponenciacio-ordinal}
Siguin \(\alpha=\operatorname{tipus}(A)\) i \(\beta=\operatorname{tipus}(B)\).
Si \(\alpha>0\), definim
\[
\alpha^{\beta}:=\operatorname{tipus}\bigl(E(B,A),<\bigr).
\]
Si \(\alpha=0\), és a dir, si \(A\) és inicial, no hi ha mínim \(0_A\) i
\(E(B,A)\) no es pot formar; definim aleshores
\[
0^{\beta}:=\operatorname{tipus}\bigl(A^{B}\bigr)
=\begin{cases}1&\text{si }\beta=0,\\0&\text{si }\beta>0,\end{cases}
\]
on \(A^{B}\), que té com a molt un punt, porta l'únic ordre possible. Els
dos valors es justifiquen dins la teoria: si \(B\) és inicial,
\(A^{B}\cong\mathbf1\) (capítol~\ref{cap:exponencials}); si \(B\) no és
inicial, un punt de \(A^{B}\) seria una aplicació \(B\to\mathbf0\), que per
l'estrictesa de \(\mathbf0\) (proposició~\ref{prop:inicial-estricte}) faria
\(B\) inicial, de manera que \(A^{B}\) no té cap punt i és inicial (lema
del buit).
\end{definicio}

\begin{teorema}[L'exponenciació està ben definida i satisfà la recursió]\label{teo:exponenciacio-recursio}
Sigui \(\alpha>0\). Per a tot bon ordre \(B\):
\begin{enumerate}
\item \(E(B,A)\) és un bon ordre, i el seu tipus no depèn dels representants
d'\(\alpha\) i \(\beta\);
\item si \(B'\hookrightarrow B\) és una part inicial, l'extensió per \(0_A\)
fora de \(B'\) identifica \(E(B',A)\) amb una part inicial de \(E(B,A)\);
\item \(\alpha^{0}=1\), \(\alpha^{\beta+1}=\alpha^{\beta}\cdot\alpha\), i, si
\(\lambda\) és límit, \(\alpha^{\lambda}=\sup_{\beta<\lambda}\alpha^{\beta}\),
on el suprem és el de la proposició~\ref{prop:suprem}.
\end{enumerate}
\emph{Axiomes usats:} exponencials, finitud (capítol~\ref{cap:finits}),
bons ordres, inducció transfinita; cap elecció; Nivell~I.
\end{teorema}

\begin{prova}
\emph{(2)} Sigui \(\iota\colon E(B',A)\to E(B,A)\) l'extensió per \(0_A\).
És estrictament creixent (els punts on difereixen dues extensions són els
mateixos). La imatge és inicial: si \(g<\iota(f)\) i \(b\) és el màxim punt
on difereixen, aleshores \(b\in B'\) ---si no, \(\iota(f)(b)=0_A\) i
\(g(b)<0_A\) seria impossible---, i per sobre de \(b\) les dues funcions
coincideixen, de manera que \(g\) també s'anul·la fora de \(B'\) (que és
inicial i conté \(b\), de manera que tot punt fora de \(B'\) és \(>b\)).

\emph{(3), cas zero i successor.} Si \(B\) és buit, \(A^{B}\) té un sol
punt, de suport buit: \(\alpha^0=1\). Si \(B=C+\mathbf1\) amb punt final
\(\ast\), l'aplicació \(f\mapsto(f|_C,f(\ast))\) és una bijecció
\(E(B,A)\cong E(C,A)\times A\) (una funció de suport finit sobre \(B\) és
una sobre \(C\) més un valor a \(\ast\)), i l'ordre antilexicogràfic, que
mira primer el punt \(\ast\), es correspon exactament amb l'ordre del
producte ordinal, on mana la coordenada d'\(A\)
(definició~\ref{def:producte-ordinal}): \(\alpha^{\beta+1}=\alpha^{\beta}\cdot\alpha\).

\emph{(3), cas límit.} Si \(B\) no té màxim, tota \(f\in E(B,A)\) té suport
finit, acotat per algun punt \(b\) de \(B\) (si el suport és no buit, el seu
màxim té un posterior, perquè \(B\) no té màxim); per tant
\(E(B,A)=\bigcup_{b\in B}E(B_{<b},A)\), una unió de parts inicials
encadenades, per~(2). Pel mateix argument de la proposició~\ref{prop:suprem}
(isomorfismes únics compatibles sobre una reunió de parts inicials), el tipus
de la reunió és el suprem dels tipus \(\alpha^{\operatorname{tipus}(B_{<b})}\),
és a dir, \(\sup_{\beta<\lambda}\alpha^{\beta}\).

\emph{(1)} Per inducció transfinita sobre els segments de \(B\) i \(B\)
mateix (teorema~\ref{teo:induccio-transfinita}): el cas buit és trivial; el
successor és un producte ordinal de bons ordres, que és un bon ordre
(proposició~\ref{prop:producte-bon-ordre}); i al límit, una reunió
encadenada de parts inicials ben ordenades és ben ordenada (el mínim d'una
part no buida es troba ja dins de qualsevol membre que la talli, com a la
prova del lema de Zorn). La independència dels representants: isomorfismes
d'ordre \(A\cong A'\) i \(B\cong B'\) indueixen, per composició, un
isomorfisme d'ordre \(E(B,A)\cong E(B',A')\) que respecta el mínim i el
suport.
\end{prova}

\begin{observacio}[Per què no definir-la directament per recursió]\label{obs:exponenciacio-acotacio}
La presentació habitual defineix \(\alpha^{\beta}\) per recursió transfinita
sobre l'exponent, amb un suprem als límits. Aquí això tindria dos
inconvenients. Primer, una recursió necessita un codomini fix i cada suprem
una família representada dins un bon ordre (observació~\ref{obs:suprem-acotat});
caldria exhibir-los. La construcció per funcions de suport finit els
proporciona d'entrada: totes les potències \(\alpha^{\beta}\) amb
\(\beta\le\lambda\) viuen, per~(2), com a parts inicials d'un sol bon ordre
\(E(L,A)\), sense cap cota externa ni cap apel·lació a SR. Segon, per a la base \(0\) les equacions del zero i del successor
continuen valent (\(0^{0}=1\) i \(0^{\beta+1}=0^{\beta}\cdot0=0\)), però la
del límit \emph{no}: amb la clàusula límit literal,
\(0^{\omega}=\sup\{0^0,0^1,0^2,\dots\}=\sup\{1,0,0,\dots\}=1\), mentre que
la definició dona \(0^{\omega}=0\), d'acord amb el fet que no hi ha cap
aplicació \(\N\to\mathbf0\). La raó és que la família \(1,0,0,\dots\) no és
creixent, i un suprem només recull bé una família creixent. Per això la
base \(0\) es defineix a part, i la clàusula límit del
teorema~\ref{teo:exponenciacio-recursio} només s'afirma per a
\(\alpha>0\), on la família \((\alpha^\beta)_\beta\) és creixent
(proposició~\ref{prop:monotonia-ordinal}). Una exponenciació amb bases i exponents variant en famílies
sense un conjunt ambient comú, en canvi, sí que demanaria SR
(observació~\ref{obs:sr}).
\end{observacio}

\begin{proposicio}[Càlculs bàsics amb l'exponenciació]\label{prop:exponenciacio-calculs}
Amb la definició~\ref{def:exponenciacio-ordinal} i el teorema~\ref{teo:exponenciacio-recursio}, per a \(\alpha>0\):
\[
\alpha^{1}=\alpha,
\qquad
\alpha^{2}=\alpha\cdot\alpha,
\qquad
2^{\omega}=\omega,
\qquad
\omega^{2}=\omega\cdot\omega .
\]
En particular \(2^{\omega}=\omega\neq\mathfrak c\): l'exponenciació
\emph{ordinal} \(2^{\omega}\) és numerable, ben diferent de l'exponenciació
\emph{cardinal} \(2^{\aleph_0}\), que és el continu.
\end{proposicio}

\begin{prova}
\(\alpha^{1}=\alpha^{0+1}=\alpha^{0}\cdot\alpha=1\cdot\alpha=\alpha\), i
\(\alpha^{2}=\alpha^{1}\cdot\alpha=\alpha\cdot\alpha\). Per a \(2^{\omega}\):
és el suprem dels \(2^{n}\) amb \(n<\omega\), tots finits, i el suprem d'una
successió creixent de tipus finits no acotada és \(\omega\); per tant
\(2^{\omega}=\omega\). I \(\omega^{2}=\omega^{1}\cdot\omega=\omega\cdot\omega\).
La comparació amb la cardinal se subratlla perquè és el nucli de la distinció
del capítol següent: \(2^{\omega}=\omega\) com a \emph{tipus ordinal} (una
longitud numerable), mentre que com a \emph{cardinal} \(2^{\aleph_0}\) és no
numerable (teorema~\ref{teo:cantor}); les dues operacions, escrites igual,
signifiquen coses distintes.
\end{prova}

\section{Monotonia i ordre contra mida}\label{sec:monotonia-ordre-mida}

\begin{lema}[Les immersions d'ordre comparen tipus]\label{lema:immersio-ordre}
Siguin \(W\) i \(V\) bons ordres i \(f\colon W\to V\) una aplicació
estrictament creixent. Aleshores \(\operatorname{tipus}(W)\le\operatorname{tipus}(V)\).
\end{lema}

\begin{prova}
Per la comparabilitat (teorema~\ref{teo:comparabilitat}), si no fos
\(\operatorname{tipus}(W)\le\operatorname{tipus}(V)\) tindríem
\(V\cong W_{<w}\) per a algun \(w\); compost amb \(f\), donaria una
aplicació estrictament creixent \(W\to W_{<w}\hookrightarrow W\) amb valor a
\(w\) estrictament menor que \(w\), contra la rigidesa
(lema~\ref{lema:rigidesa}).
\end{prova}

\begin{proposicio}[Monotonia]\label{prop:monotonia-ordinal}
Les tres operacions són \emph{estrictament creixents per la dreta} i
\emph{no decreixents per l'esquerra}. En particular, si \(\beta<\gamma\):
\[
\alpha+\beta<\alpha+\gamma,
\qquad
\alpha\cdot\beta<\alpha\cdot\gamma\ (\alpha>0),
\qquad
\alpha^{\beta}<\alpha^{\gamma}\ (\alpha>1),
\]
mentre que per l'esquerra només val \(\le\): de \(\beta<\gamma\) se'n
segueix \(\beta+\alpha\le\gamma+\alpha\), i pot fallar l'estrictesa
(\(1<2\) però \(1+\omega=2+\omega=\omega\)).
\end{proposicio}

\begin{prova}
Siguin \(A,B,C\) representants d'\(\alpha,\beta,\gamma\), amb
\(B\cong C_{<c}\) per a un punt \(c\) de \(C\) (perquè \(\beta<\gamma\)).

\emph{Suma per la dreta.} \(A+C_{<c}\) és el segment de \(A+C\) determinat
pel punt \(c\) del segon sumand, que és propi; d'on
\(\alpha+\beta<\alpha+\gamma\).

\emph{Producte per la dreta.} Si \(\alpha>0\), sigui \(a_0\) el mínim d'\(A\).
En l'ordre antilexicogràfic, \(A\times C_{<c}\) és el segment de
\(A\times C\) determinat per \((a_0,c)\), que és propi; d'on
\(\alpha\cdot\beta<\alpha\cdot\gamma\).

\emph{Exponenciació per la dreta.} Si \(\alpha>1\), el representant \(A\)
té, a més del mínim \(0_A\), algun punt \(a_1\neq0_A\). Per la invariància
(teorema~\ref{teo:exponenciacio-recursio}(1)),
\(\alpha^{\beta}=\operatorname{tipus}(E(C_{<c},A))\), i per
l'apartat~(2) d'aquell teorema l'extensió per \(0_A\) identifica
\(E(C_{<c},A)\) amb una part inicial de \(E(C,A)\). Aquesta part és
\emph{pròpia}: l'aplicació \(g\colon C\to A\) que val \(a_1\) al punt \(c\) i
\(0_A\) a la resta (definició per casos sobre el singletó \(\{c\}\) i el seu
complement) té suport \(\{c\}\), finit, i no és l'extensió de cap funció
sobre \(C_{<c}\), perquè \(g(c)\neq0_A\). Per la
proposició~\ref{prop:parts-inicials}, \(E(C_{<c},A)\) és isomorf a un segment
propi de \(E(C,A)\), és a dir, \(\alpha^{\beta}<\alpha^{\gamma}\). No cal cap
inducció sobre els tipus, que no formen cap conjunt: l'argument passa
íntegrament dins el bon ordre fix \(C\).

\emph{Per l'esquerra.} Si \(\beta\le\gamma\), hi ha una immersió estrictament
creixent \(B\to C\) (sobre un segment); en resulten immersions estrictament
creixents \(B+A\to C+A\) (per a \(\beta+\alpha\le\gamma+\alpha\)) i
\(B\times A\to C\times A\) en l'ordre antilexicogràfic (per a
\(\beta\cdot\alpha\le\gamma\cdot\alpha\)),
i pel lema~\ref{lema:immersio-ordre} els tipus es comparen. Per a
l'exponenciació, sigui \(D\) un representant de \(\delta\); volem
\(\beta^{\delta}\le\gamma^{\delta}\). Si \(\beta=0\), \(\beta^{\delta}\) val
\(1\) o \(0\) segons si \(\delta=0\) o no
(definició~\ref{def:exponenciacio-ordinal}), i en el primer cas també
\(\gamma^{0}=1\); la desigualtat és immediata. Si \(\beta>0\), la immersió
\(j\colon B\to C\) és un isomorfisme d'ordre sobre una part inicial no
inicial de \(C\), que conté el mínim de \(C\); per tant \(j(0_B)=0_C\). La
postcomposició \(f\mapsto j\circ f\) envia \(E(D,B)\) dins \(E(D,C)\): conserva
el suport, perquè \(j\) és injectiva i \(j(0_B)=0_C\), i és estrictament
creixent, perquè \(f\) i \(f'\) difereixen exactament als mateixos punts que
\(j\circ f\) i \(j\circ f'\), i al màxim d'aquests punts \(j\) conserva la
desigualtat. Pel lema~\ref{lema:immersio-ordre},
\(\beta^{\delta}\le\gamma^{\delta}\). L'exemple \(1+\omega=2+\omega=\omega\)
(proposició~\ref{prop:suma-no-commutativa}) mostra que per l'esquerra
l'estrictesa pot fallar.
\end{prova}

\begin{observacio}[Ordre contra mida: la lliçó de l'aritmètica ordinal]\label{obs:ordre-vs-mida}
L'aritmètica ordinal \emph{recorda l'ordre}; la cardinal, que veurem tot
seguit, l'\emph{oblidarà}. Els càlculs d'aquest capítol ho fan tangible.
Ordinalment, \(1+\omega=\omega\) però \(\omega+1>\omega\);
\(2\cdot\omega=\omega\) però \(\omega\cdot2>\omega\); \(2^{\omega}=\omega\).
Els tipus \(\omega\), \(\omega+1\), \(\omega\cdot2\), \(\omega^{2}\) són
\emph{diferents} dos a dos com a formes d'ordre; en canvi \(2^{\omega}=\omega\)
(exponenciació \emph{ordinal}!), de manera que \(2^{\omega}\) no és un tipus
nou sinó el mateix \(\omega\). Tots ells són a més \emph{equipotents}: tots
numerables, tots de cardinal \(\aleph_0\). Com a \emph{mides} són el mateix. Aquesta és la raó de mantenir separades les dues
aritmètiques: la ordinal distingeix \(\omega+1\) de \(\omega\), cosa que la
cardinal no pot fer, perquè oblidar l'ordre és precisament identificar tots
els bons ordres numerables. El capítol següent construeix l'aritmètica
cardinal, on \enquote{\(+\)}, \enquote{\(\cdot\)} i \enquote{\(\,\hat{}\,\)}
signifiquen suma disjunta, producte i exponencial de conjunts \emph{sense}
ordre, i on \(2^{\aleph_0}\neq\aleph_0\) per Cantor.
\end{observacio}

\begin{resumcap}
\apartat{Idees centrals}
\begin{enumerate}
  \item Siguin \(A\) i \(B\) representants dels tipus \(\alpha\) i \(\beta\). La suma \(\alpha+\beta\) és el tipus de la concatenació \(A+B\) (tot \(A\) abans de tot \(B\)). El producte \(\alpha\cdot\beta\) és el tipus de \(A\times B\) amb l'ordre antilexicogràfic: una còpia d'\(A\) per cada punt de \(B\). Totes dues operacions no depenen dels representants i són associatives.
  \item No són commutatives: \(1+\omega=\omega<\omega+1\), \(2\cdot\omega=\omega<\omega\cdot2\). El producte distribueix només per l'esquerra.
  \item Per a \(\alpha>0\), l'exponenciació \(\alpha^\beta\) és el tipus del conjunt \(E(B,A)\) de les aplicacions de suport finit \(B\to A\), amb l'ordre antilexicogràfic; satisfà \(\alpha^0=1\), \(\alpha^{\beta+1}=\alpha^\beta\cdot\alpha\) i el suprem als límits, sense cap cota externa. Per a \(\alpha=0\), \(0^\beta\) és el tipus de \(A^B\): \(1\) si \(\beta=0\) i \(0\) altrament; la clàusula límit no val per a aquesta base. Així \(2^\omega=\omega\) ordinalment.
  \item Una immersió estrictament creixent compara tipus. Les tres operacions són estrictament creixents per la dreta i no decreixents per l'esquerra.
  \item L'aritmètica ordinal recorda l'ordre; la cardinal, que ve tot seguit, l'oblida: \(\omega\), \(\omega+1\), \(\omega\cdot2\), \(\omega^2\) són tipus diferents i la mateixa mida.
\end{enumerate}
\apartat{Errors típics} Invertir la convenció del producte (\(\omega\cdot2\) són dues còpies de \(\omega\), no \(\omega\) còpies de \(2\)); pensar que \(\omega+1=1+\omega\); confondre l'exponenciació ordinal \(2^\omega=\omega\) amb la cardinal \(2^{\aleph_0}\); aplicar la clàusula límit a la base \(0\); i creure que la monotonia per l'esquerra és estricta.
\apartat{Lectura recomanada} \cite[§§19--21]{halmos} per a l'aritmètica ordinal en la teoria clàssica; \cite[cap.~2]{jech2003} per a la suma, el producte i l'exponenciació ordinals, amb la forma normal de Cantor; i \cite[cap.~I.9]{kunen2011} per a la definició recursiva de les operacions i les seves propietats.
\end{resumcap}
